\clearpage
\fancyhead[L]{\small BOUVARD. JUPITER UND SATURN.}
\SourceLeaf{513}
\begin{center}\small Göttingische gelehrte Anzeigen. Stück 177. S. 1761..1764. 5. November 1808.\\[3pt]\rule{.72\linewidth}{.3pt}\end{center}

Bei Courcier. \emph{Tables astronomiques publiées par le bureau des longitudes de France. Nouvelles tables de Jupiter et de Saturne calculées d'après la théorie de M. }\textsc{Laplace}, \emph{et suivant la division decimale de l'angle droit; par M.} \textsc{Bouvard}. 1808. Im Quartformat. Paris.

Diese Tafeln sind als Fortsetzung der Sonnen- und Mondtafeln anzusehen, die 1806 herausgegeben und bereits im 95. Stück dieser Blätter von 1807 besprochen wurden. Sie sind die kostbare Frucht der Untersuchungen, die \textsc{Laplace} mehr als 20 Jahre lang fortgesetzt und immer weiter verfeinert hat. Sie betreffen die Störungen, durch die die Bewegungen der beiden größten Planeten unseres Sonnensystems so verwickelt werden. Es ist bekannt, mit welch glücklichem Erfolg dieser große Geometer hier wie in allen anderen Gebieten der physikalischen Astronomie, die er erforscht hat, alle Schwierigkeiten überwunden und die widerspenstigen Bewegungen der Rechnung unterworfen hat. Ihm verdanken wir es, dass wir diese Bewegung jetzt fast ebenso genau berechnen wie beobachten können. Die auf den ersten Untersuchungen von \textsc{Laplace} beruhenden Jupiter- und Saturntafeln, die \textsc{Delambre} berechnet hatte, wichen schon selten um eine halbe Minute vom Himmel ab. Seitdem hat \textsc{Laplace} die Berechnung der Störungen noch viel weiter geführt. Unter Berücksichtigung dieser Rechnungen hat \textsc{Bouvard} aus einer äußerst sorgfältigen Auswertung aller seit einem halben Jahrhundert beobachteten Oppositionen die elliptischen Elemente neu bestimmt. Das hat eine dreimal so große Übereinstimmung ergeben. Die letzten Ergebnisse dieser Arbeit, nämlich die in Formeln gefassten Bewegungen von Jupiter und Saturn, sind bereits im 8. Teil der Mécanique céleste, S. 337, veröffentlicht worden. Im vorliegenden Werk finden wir sie nun in Tafeln. Über deren Einrichtung haben wir nur noch wenig hinzuzufügen, da sie in den meisten Punkten mit der aus den Sonnentafeln von \textsc{Delambre} bekannten Einrichtung übereinstimmt. Man zählt die Zeiten ab Mitternacht statt ab Mittag. Die von der Sonnennähe aus gemessenen Anomalien und die Gleichungen wurden alle positiv gemacht, indem man zu jeder eine konstante Größe addierte und die Summe aller dieser Zusätze von der Mittelpunktsgleichung (die, wo nötig, durch Hinzufügen von 400 Dezimalgraden positiv gemacht wurde), vom elliptischen Radiusvektor und vom Abstand vom Nordpol, der an die Stelle der Breite tritt, wieder abzog. Alle diese Festlegungen haben seit Langem den Beifall der Astronomen gefunden, abgesehen von der neuen Art der Zeitrechnung, die man ohne irgendeinen erkennbaren Nutzen eingeführt hat. Mit der Dezimaleinteilung des Quadranten in diesen 'Tafeln will man versuchen, die Astronomen nach und nach daran zu gewöhnen. Solange jedoch zugleich Sonnen- und Sinustafeln nach der alten Einteilung gebraucht werden müssen, ist dies eher eine Unbequemlichkeit als eine Erleichterung.

Die den Tafeln vorangestellte Einleitung beschreibt zunächst die einzelnen Tafeln. Darauf folgen die Formeln, nach denen sie berechnet wurden. Sie unterscheiden sich von den an der genannten Stelle der Mécanique céleste angegebenen nur dadurch, dass man alle Epochen und veränderlichen Koeffizienten von 1750 auf 1800 reduziert hat. Schließlich folgt die vollständige Berechnung einer Jupiterbeobachtung vom 2. April 1806, bei der die Tafeln um $20^{\prime\prime}0$ in der Länge und um $19^{\prime\prime}7$ in der Breite abweichen ($6^{\prime\prime}8$ und $6^{\prime\prime}4$ nach der Sexagesimaleinteilung). Das ist der Fehler für die geozentrische Position. Anschließend wird gezeigt, wie man daraus den Fehler der heliozentrischen Position ableiten kann. Da man dabei aber den berechneten Radiusvektor als vollkommen genau ansehen muss, hätte man darauf hinweisen müssen, dass sich dieses Verfahren nur dort mit einiger Sicherheit anwen\SourceLeaf{514}den lässt, wo ein mäßiger Fehler im Radiusvektor nur einen sehr geringen Einfluss auf Länge und Breite hat, also nur bei entfernten Planeten oder in der Nähe der Opposition sowie bei mäßigen Neigungen oder in der Nähe der Knoten.

Da diese Jupiter- und Saturntafeln allein auf den neueren Beobachtungen beruhen, ist es interessant zu sehen, mit welcher Genauigkeit sie ältere Beobachtungen wiedergeben. Besonders bemerkenswert ist in dieser Hinsicht die Begegnung der beiden Planeten, die \textsc{Ibn Junis} am 31. Oktober 1007 in Kairo beobachtete. Nach der Rechnung von \textsc{Bouvard} ergibt sich hier zwischen den Tafeln und der Beobachtung ein Unterschied von $11'\,25''$ in der Länge und von $6'\,56''$ in der Breite (Sexagesimaleinteilung), der den möglichen Beobachtungsfehler kaum übersteigt. Wir bemerken hierzu noch, dass sich nach dem Erscheinen der vorliegenden Tafeln ein kleiner Fehler in der großen Gleichung sowohl für Jupiter als auch für Saturn gefunden hat. Ein Glied war mit falschem Vorzeichen angesetzt worden. Diese Änderung in der großen Gleichung macht aber zugleich eine kleine Änderung der Epochen und mittleren Bewegungen nötig. Demnach muss zu $q^{\mathrm{IV}}$ (Mécan. cél., Band 4, S. 338) noch Folgendes hinzugefügt werden
\[
\begin{aligned}
&51''98 + t.0''4156\\
&{}- (73''58 - t.0''0103)\sin(5n^{\mathrm{V}}t - 2n^{\mathrm{IV}}t + 5\varepsilon^{\mathrm{V}} - 2\varepsilon^{\mathrm{IV}})\\
&{}+ (47''63 + t.0''0287)\cos(5n^{\mathrm{V}}t - 2n^{\mathrm{IV}}t + 5\varepsilon^{\mathrm{V}} - 2\varepsilon^{\mathrm{IV}})
\end{aligned}
\]
und zu $q^{\mathrm{V}}$ (ebenda)
\[
\begin{aligned}
&-127''13 - t.1''0212\\
&{}+ (179''952 - t.0''025192)\sin(5n^{\mathrm{V}}t - 2n^{\mathrm{IV}}t + 5\varepsilon^{\mathrm{V}} - 2\varepsilon^{\mathrm{IV}})\\
&{}- (116''541 + t.0''070196)\cos(5n^{\mathrm{V}}t - 2n^{\mathrm{IV}}t + 5\varepsilon^{\mathrm{V}} - 2\varepsilon^{\mathrm{IV}})
\end{aligned}
\]
Dadurch wird der oben genannte Fehler in der Länge auf $5'\,22''$ verringert und ist damit so klein, dass man ihn ohne Bedenken der Beobachtung zuschreiben kann.

